\documentclass[../main.tex]{subfiles}
\begin{document}

\begin{center}
    \Large{\textbf{ABSTRACT}}\\
\end{center}
\vspace{1cm}
Large language models have become effective systems for text generation and reasoning, but
their internal computations remain difficult to inspect. This opacity limits the ability of
researchers to explain why a model produces a prediction, diagnose failures, and build
reliable evidence about model behavior. Mechanistic interpretability addresses this problem
by reverse-engineering model computations into human-understandable mechanisms. Within that
field, circuit tracing studies the specific internal pathways that support one behavior or
one prediction. Recent circuit-tracing methods use replacement models based on
transcoders to express a model's computation as an attribution graph, whose nodes are
interpretable features and whose signed edges record linear effects leading to a
target token.

Attribution graphs provide a faithful and prompt-specific view of model computation, but
they are still too large for direct human reading. A single prompt can produce thousands of
feature nodes and many weighted edges, so existing workflows rely on manual pruning and
manual grouping by expert inspection. This thesis chooses automated attribution-graph
summarization as the central approach because it preserves the mechanistic object produced
by circuit tracing while replacing the most time-consuming human steps with explicit
objectives. The proposed method takes an attribution graph as input and returns a compact
summary graph. It first prunes nodes and edges using influence and relevance scores, then
clusters surviving feature nodes into supernodes under an objective that rewards atomicity
and causal preservation, and finally labels and assembles the resulting directed acyclic
summary graph for inspection.

The main contributions of this thesis are a formalization of faithful attribution-graph
summarization, an automated three-stage pipeline for pruning, clustering, and visualizing
such graphs, and an empirical evaluation on analogy and multi-hop reasoning prompts. The
results show that the summaries preserve graph-level faithfulness, improve structural
coherence over clustering baselines, and retain signed intervention behavior with high sign
consistency and positive edge-effect rank correlation.

\begin{flushright}
\begin{tabular}{@{}c@{}}
Student\\
\textit{(Signature and full name)}
\end{tabular}
\end{flushright}

\end{document}
